\section{Step 51 Reject a reconstructed value below the signed minimum}\label{delta-s51}

Load the full signed minimum into \texttt{t7} and compare \texttt{t6} with it. Use \texttt{dec\_range} when the candidate is smaller. The canonical 33-bit mapped maximum decodes to a delta too negative for any allowed predecessor and is a useful test of this distinction.

\Needspace{10\baselineskip}
\subsection{Write and inspect}
Replace the supplied temporary branch and nop in \texttt{d\_minimum}. Keep surrounding labels and region markers.

\MARSCodeLabel{MIPS assembly}
\begin{minted}[samepage=true]{mars-idea}
    dli $t7, -2147483648
    slt $t9, $t6, $t7
    bne $t9, $zero, dec_range
    nop
\end{minted}
Use \texttt{s51-1}, \texttt{s51-2} from the companion construction inputs. Pause at \texttt{codec\_done} before caller result storage.

\subsection{Expected observation}
Canonical ff ff ff ff 1f returns range status three rather than malformed status two.

\Needspace{11\baselineskip}
\subsection{Check your understanding}
Why is that five-byte sequence not malformed?

\begin{itemize}
\item \textbf{A} Its width, continuation and minimality are valid.
\item \textbf{B} Every five-byte sequence is valid.
\item \textbf{C} The output capacity is large.
\end{itemize}
